\documentclass[11pt,a4paper]{article}

\usepackage{fontspec}
\IfFontExistsTF{Times New Roman}{\setmainfont{Times New Roman}}{\setmainfont{Tinos}}
\IfFontExistsTF{Consolas}{\setmonofont{Consolas}[Scale=0.92]}{\setmonofont{Liberation Mono}[Scale=0.92]}
\usepackage[english]{babel}
\usepackage{geometry}
\usepackage{xcolor}
\usepackage{enumitem}
\usepackage{listings}
\usepackage{booktabs}
\usepackage{array}
\usepackage{amsmath,amssymb}
\usepackage{titlesec}
\usepackage{fancyhdr}
\usepackage{parskip}
\usepackage{tcolorbox}
\tcbuselibrary{skins}
\usepackage{tabularx}
\usepackage{lastpage}
\usepackage{hyperref}

\geometry{top=1.5cm,bottom=1.5cm,left=1.7cm,right=1.7cm}

\definecolor{primaryblue}{RGB}{18,65,115}
\definecolor{secondaryteal}{RGB}{14,102,110}
\definecolor{SoftBlue}{RGB}{232,242,250}
\definecolor{CodeBg}{RGB}{248,250,252}
\definecolor{DarkGray}{RGB}{55,55,55}
\definecolor{boxbg}{RGB}{248,250,252}
\definecolor{borderblue}{RGB}{150,185,215}

\setlength{\parindent}{0pt}
\setlength{\parskip}{0pt}
\setlist[itemize]{leftmargin=1.3em,itemsep=1pt,topsep=1pt,parsep=0pt,partopsep=0pt}
\setlist[enumerate]{leftmargin=1.5em,itemsep=1.5pt,topsep=1pt,parsep=0pt,partopsep=0pt}
\setlist[enumerate,2]{topsep=0.5pt,itemsep=0pt,parsep=0pt,partopsep=0pt}

\lstset{
  language=Python,
  basicstyle=\ttfamily\scriptsize,
  keywordstyle=\color{blue!75!black}\bfseries,
  stringstyle=\color{red!65!black},
  commentstyle=\color{primaryblue!80!black}\itshape,
  showstringspaces=false,
  columns=fullflexible,
  keepspaces=true,
  breaklines=true,
  frame=single,
  rulecolor=\color{primaryblue!45},
  backgroundcolor=\color{CodeBg},
  xleftmargin=0.2em,
  xrightmargin=0.2em,
  aboveskip=0.2em,
  belowskip=0.2em,
  tabsize=4
}

\pagestyle{fancy}
\fancyhf{}
\setlength{\headheight}{14pt}
\fancyhead[L]{\footnotesize\textbf{DSAI1003 -- Fundamental Programming Concepts in Python}}
\fancyhead[R]{\footnotesize\textbf{Lists \& Tuples Exam -- 60 Minutes}}
\fancyfoot[L]{\scriptsize Faculty of Data Science and Artificial Intelligence -- National Economics University}
\fancyfoot[R]{\footnotesize Page \thepage\ of 6}
\renewcommand{\headrulewidth}{0.4pt}
\renewcommand{\footrulewidth}{0.3pt}

\titleformat{\section}{\large\bfseries\color{primaryblue}}{}{0pt}{}
\titleformat{\subsection}{\normalsize\bfseries\color{primaryblue}}{}{0pt}{}
\titlespacing*{\section}{0pt}{3.5pt}{1.5pt}
\titlespacing*{\subsection}{0pt}{2.5pt}{1pt}

\newtcolorbox{headerbox}{
  colback=white,
  colframe=primaryblue,
  boxrule=0.7pt,
  arc=1.5mm,
  left=2.5mm,
  right=2.5mm,
  top=1.2mm,
  bottom=1.2mm
}

\newtcolorbox{answerbox}[2][]{
  enhanced,
  colback=white,
  colframe=DarkGray!45,
  boxrule=0.5pt,
  arc=1mm,
  left=2mm,
  right=2mm,
  top=1.2mm,
  bottom=1.2mm,
  title=\footnotesize\color{DarkGray}\textbf{#2},
  #1
}

\hypersetup{
  colorlinks=true,
  linkcolor=primaryblue,
  urlcolor=primaryblue,
  pdftitle={DSAI1003 - Lists, Tuples and Sequential Data Structures in Python Exam},
  pdfauthor={Faculty of Data Science and Artificial Intelligence - National Economics University (NEU)}
}

\begin{document}

% =========================================================================
% PAGE 1: OFFICIAL HEADER + ANSWER SHEET + SECTION A (Q1 - Q5)
% =========================================================================

\noindent
\begin{minipage}[t]{0.64\textwidth}
  \raggedright\small
  \textbf{NATIONAL ECONOMICS UNIVERSITY}\\
  \textbf{COLLEGE OF TECHNOLOGY}\\
  \textbf{FACULTY OF DATA SCIENCE \& ARTIFICIAL INTELLIGENCE (FDA)}
\end{minipage}%
\hfill
\begin{minipage}[t]{0.34\textwidth}
  \raggedleft\footnotesize
  \textbf{OFFICIAL PROGRESS EXAMINATION}\\
  Academic Year: 2026--2027\\
  Semester: Fall 2026 \quad Course Code: \textbf{DSAI1003}
\end{minipage}

\vspace{0.15em}

\begin{center}
  {\Large\bfseries\color{primaryblue} FUNDAMENTAL PROGRAMMING CONCEPTS IN PYTHON}\\[0.08em]
  {\normalsize\bfseries PROGRESS EXAM: LISTS, TUPLES \& SEQUENTIAL DATA STRUCTURES}\\[0.08em]
  {\normalsize\textbf{Paper Code:} \rule{1.6cm}{0.4pt} \quad | \quad \textbf{Duration: 60 Minutes} (Excluding paper distribution)}
\end{center}

\vspace{0.05em}

\begin{headerbox}
\begin{tabular*}{\textwidth}{@{\extracolsep{\fill}}ll@{}}
  \textbf{Candidate Full Name:} \dotfill & \textbf{Student ID (MSSV):} \dotfill \\[0.2em]
  \textbf{Class / Cohort:} \dotfill & \textbf{Exam Room / Seat No.:} \dotfill
\end{tabular*}

\vspace{0.2em}
\renewcommand{\arraystretch}{1.05}
\setlength{\tabcolsep}{4.5pt}
\centering
\begin{tabular}{|l|c|c|c|c||c|c|}
\hline
\textbf{Section} & \textbf{A} & \textbf{B} & \textbf{C} & \textbf{D} & \textbf{Total (100)} & \textbf{Examiner Signature} \\
\hline
\textbf{Max Score} & 20 & 25 & 25 & 30 & \textbf{100} & \\[0.3em]
\cline{1-6}
\textbf{Score} & & & & & & \\[0.4em]
\hline
\end{tabular}

\vspace{0.15em}
\raggedright
\scriptsize
\textbf{Instructions:} Closed-book examination. All electronic devices and internet access are prohibited. Answer all questions directly in the designated spaces. In programming tasks, apply indexing, slicing, mutability rules, list methods (\texttt{append}, \texttt{insert}, \texttt{pop}, \texttt{remove}), linear search, 2D matrices, and tuple packing/unpacking. Do not detach pages.
\end{headerbox}

\vspace{1mm}
\noindent
\textbf{MULTIPLE-CHOICE ANSWER SHEET} \textit{(Write A, B, C, or D in the corresponding box)}:
\begin{center}
\setlength{\tabcolsep}{5pt}
\renewcommand{\arraystretch}{1.0}
\begin{tabular}{|c|c|c|c|c|c|c|c|c|c|c|}
\hline
\textbf{Question} & \textbf{1} & \textbf{2} & \textbf{3} & \textbf{4} & \textbf{5} & \textbf{6} & \textbf{7} & \textbf{8} & \textbf{9} & \textbf{10} \\
\hline
\textbf{Answer} & & & & & & & & & & \\[1.8mm]
\hline
\end{tabular}
\end{center}
\vspace{-2.5mm}
\section*{Section A: Sequential Collections, Lists \& Tuples Concepts \hfill \normalsize\normalfont(20 points)}
\textit{Answer all 10 questions. Each correct answer carries 2 points. Suggested time: 12 minutes.}

\begin{enumerate}[label=\textbf{Q\arabic*.},leftmargin=2em,itemsep=1.5pt,topsep=1pt]

\item Given the list \texttt{nums = [10, 20, 30, 40, 50]}, what does the slicing expression \texttt{nums[1:4:2]} return?
\begin{enumerate}[label=\Alph*.,leftmargin=2em]
  \item \texttt{[20, 40]}
  \item \texttt{[20, 30, 40]}
  \item \texttt{[10, 30]}
  \item \texttt{[20, 50]}
\end{enumerate}

\item What is the fundamental difference between a Python \textbf{list} and a \textbf{tuple}?
\begin{enumerate}[label=\Alph*.,leftmargin=2em]
  \item Lists can store only numbers, whereas tuples can store arbitrary heterogeneous data types.
  \item Lists are mutable (elements can be modified, appended, or removed in-place), whereas tuples are strictly immutable.
  \item Tuples do not support zero-based indexing or negative index positions.
  \item Lists are delimited by parentheses \texttt{()}, whereas tuples require square brackets \texttt{[]}.
\end{enumerate}

\item Given \texttt{a = [1, 2, 3]} and the assignment \texttt{b = a}, what happens when \texttt{b.append(4)} is executed?
\begin{enumerate}[label=\Alph*.,leftmargin=2em]
  \item \texttt{b} becomes \texttt{[1, 2, 3, 4]}, while \texttt{a} remains \texttt{[1, 2, 3]} because \texttt{b} is an independent copy.
  \item Both \texttt{a} and \texttt{b} reference the same underlying list object, so both reflect \texttt{[1, 2, 3, 4]}.
  \item A \texttt{TypeError} is raised because variable aliasing is prohibited for mutable structures in Python.
  \item \texttt{a} is automatically converted into an immutable tuple to protect original values.
\end{enumerate}

\item How does the method \texttt{nums.append([4, 5])} differ from \texttt{nums.extend([4, 5])} when applied to \texttt{nums = [1, 2, 3]}?
\begin{enumerate}[label=\Alph*.,leftmargin=2em]
  \item \texttt{append} adds the sublist as a single nested element \texttt{[1, 2, 3, [4, 5]]}; \texttt{extend} appends elements individually to form \texttt{[1, 2, 3, 4, 5]}.
  \item Both methods produce identical flat lists of length 5.
  \item \texttt{extend} raises an exception because it operates exclusively on strings, not lists.
  \item \texttt{append} inserts elements at index 0, whereas \texttt{extend} inserts elements at the tail.
\end{enumerate}

\item What is the effect of executing \texttt{item = lst.pop(1)} on \texttt{lst = ["a", "b", "c", "d"]}?
\begin{enumerate}[label=\Alph*.,leftmargin=2em]
  \item It searches for the value \texttt{1} and deletes it; \texttt{item} receives \texttt{None}.
  \item It removes the element at index 1 (\texttt{"b"}), returns \texttt{"b"} to \texttt{item}, and shifts subsequent elements left.
  \item It removes the first element (\texttt{"a"}) and assigns the remaining sublist to \texttt{item}.
  \item It removes elements from index 1 through the end of the list.
\end{enumerate}

\end{enumerate}

\newpage

% ==============================================================================
% PAGE 2: SECTION A (Q6 - Q10) + SECTION B Q11
% ==============================================================================

\begin{enumerate}[label=\textbf{Q\arabic*.},leftmargin=2em,itemsep=1.5pt,topsep=1pt,resume]

\item What happens when executing \texttt{t = (10, 20, 30); t[0] = 99}?
\begin{enumerate}[label=\Alph*.,leftmargin=2em]
  \item \texttt{t} is modified to \texttt{(99, 20, 30)}.
  \item A new tuple is implicitly created and rebound to \texttt{t}.
  \item A \texttt{TypeError} is raised because tuples do not support item assignment.
  \item \texttt{t} is silently converted into a Python list.
\end{enumerate}

\item What values are assigned to \texttt{x}, \texttt{y}, and \texttt{z} by the statement \texttt{x, y, *z = (1, 2, 3, 4, 5)}?
\begin{enumerate}[label=\Alph*.,leftmargin=2em]
  \item \texttt{x = 1}, \texttt{y = 2}, and \texttt{z = [3, 4, 5]}
  \item \texttt{x = (1, 2)}, \texttt{y = 3}, and \texttt{z = (4, 5)}
  \item A \texttt{ValueError} is raised because tuple unpacking requires exact matching lengths.
  \item \texttt{x = 1}, \texttt{y = [2, 3]}, and \texttt{z = 5}
\end{enumerate}

\item What occurs when calling \texttt{names.index("Charlie")} if \texttt{"Charlie"} is \textbf{not} present in \texttt{names}?
\begin{enumerate}[label=\Alph*.,leftmargin=2em]
  \item The method returns \texttt{-1} to indicate absence.
  \item The method returns \texttt{None}.
  \item A \texttt{ValueError} exception is raised.
  \item The string is automatically appended to the end of the list.
\end{enumerate}

\item In a 2D matrix represented as a list of lists \texttt{grid = [[1, 2, 3], [4, 5, 6], [7, 8, 9]]}, how do you access the value \texttt{6}?
\begin{enumerate}[label=\Alph*.,leftmargin=2em]
  \item \texttt{grid[2][1]}
  \item \texttt{grid[1][2]}
  \item \texttt{grid[2, 1]}
  \item \texttt{grid[6]}
\end{enumerate}

\item When passing a list to a function in Python, how are modifications inside the function handled?
\begin{enumerate}[label=\Alph*.,leftmargin=2em]
  \item Lists are passed strictly by value; in-place modifications never affect the caller's list.
  \item A shallow copy is automatically made upon entering the function body.
  \item The list reference is passed; modifying elements in-place modifies the original list object in the caller.
  \item In-place mutations raise an \texttt{UnboundLocalError} unless declared with the \texttt{global} keyword.
\end{enumerate}

\end{enumerate}

\vspace{1mm}
\section*{Section B: Output Tracing, Logic Diagnosis \& Error Correction \hfill \normalsize\normalfont(25 points)}
\textit{Answer all questions Q11--Q15 directly in the designated spaces. Suggested time: 15 minutes.}

\subsection*{Q11. Exact Output Tracing with Slicing, In-Place Mutation \& Aliasing (5 points)}
Trace the following program carefully and write its exact console output:

\begin{lstlisting}
data = [10, 25, 40, 55, 70]
sub = data[1:4]
sub.append(99)
data.insert(2, 33)
popped = data.pop()
alias = data
alias[0] = 100
print(f"Data: {data}")
print(f"Sub: {sub}")
print(f"Popped: {popped}")
\end{lstlisting}

\begin{answerbox}[height=3.4cm]{Output Box for Q11}
\end{answerbox}

\newpage

% ==============================================================================
% PAGE 3: SECTION B, Q12--Q15
% ==============================================================================

\subsection*{Q12. 2D Matrix / Grid Processing Hand-Trace Table (5 points)}
Trace the following nested matrix manipulation. Fill in the state trace table for each row iteration:

\begin{minipage}[t]{0.45\textwidth}
\begin{lstlisting}
matrix = [
    [2, 5, 1],
    [4, 3, 6],
    [7, 0, 8]
]
diag_sum = 0
row_maxima = []
for r in range(len(matrix)):
    m = matrix[r][0]
    for c in range(len(matrix[r])):
        if matrix[r][c] > m:
            m = matrix[r][c]
        if r == c:
            diag_sum += matrix[r][c]
    row_maxima.append(m)
\end{lstlisting}
\end{minipage}%
\hfill
\begin{minipage}[t]{0.52\textwidth}
\centering
\footnotesize
\setlength{\tabcolsep}{3.5pt}
\begin{tabular}{|c|c|c|c|}
\hline
\textbf{Row $r$} & \textbf{Diag Value} & \textbf{Row Max $m$} & \textbf{Running \texttt{diag\_sum}} \\
\hline
0 & & & \\[1.5mm]
\hline
1 & & & \\[1.5mm]
\hline
2 & & & \\[1.5mm]
\hline
\multicolumn{4}{|l|}{\textbf{Final \texttt{diag\_sum}:} \dotfill{} \textbf{\texttt{row\_maxima}:} \dotfill{}} \\[1.5mm]
\hline
\end{tabular}
\end{minipage}

\vspace{1mm}
\subsection*{Q13. Diagnose and Correct Three Sequential Data Defects (5 points)}
Identify the exact defect and write the corrected code snippet in the table below:

\renewcommand{\arraystretch}{1.15}
\begin{tabularx}{\textwidth}{|X|p{4.2cm}|p{6.0cm}|}
\hline
\textbf{Flawed Snippet} & \textbf{Defect Classification} & \textbf{Exact Corrected Line(s)} \\
\hline
\texttt{nums = [1, 2, 3, 4, 5]}\\
\texttt{for x in nums:}\\
\texttt{\ \ \ \ if x \% 2 == 0: nums.remove(x)}\\
(Intended: remove all even items) & & \\[3.2mm]
\hline
\texttt{point = (10, 20)}\\
\texttt{point[0] += 5}\\
(Intended: translate x by +5) & & \\[3.2mm]
\hline
\texttt{orig = [1, 2, [3, 4]]}\\
\texttt{clone = orig[:]}\\
\texttt{clone[2].append(5)}\\
(Intended: independent sublist) & & \\[3.2mm]
\hline
\end{tabularx}

\vspace{1mm}
\subsection*{Q14. List Filtering, Transformation \& Accumulation (5 points)}
Predict the exact printed output of the following list processing sequence:

\begin{lstlisting}
scores = [45, 82, 67, 91, 58, 74, 88]
passing = [s for s in scores if s >= 60]
curved = [min(100, s + 5) for s in passing]
avg_curved = sum(curved) / len(curved)
print(f"Passing Count: {len(passing)}")
print(f"Curved Scores: {curved}")
print(f"Average: {avg_curved:.2f}")
\end{lstlisting}

\begin{answerbox}[height=2.2cm]{Output Box for Q14}
\end{answerbox}

\vspace{1mm}
\subsection*{Q15. Operational Comparison: Lists versus Tuples (5 points)}
\begin{itemize}[leftmargin=2em]
  \item State two distinct operational scenarios in software design where a \textbf{tuple} is strictly preferred over a \textbf{list}: \hrulefill
  \item Why can a tuple containing only integers be used as a key in a Python dictionary, while a list of integers cannot? \hrulefill
\end{itemize}

\newpage

% ==============================================================================
% PAGE 4: SECTION C
% ==============================================================================

\section*{Section C: Program Design \& Sequential Data Processing \hfill \normalsize\normalfont(25 points)}
\textit{Answer all tasks. Suggested time: 15 minutes.}

\begin{tcolorbox}[colback=boxbg,colframe=secondaryteal,arc=1.5mm,boxrule=0.6pt,top=1.5mm,bottom=1.5mm]
\textbf{Problem Specification: Retail Sensor Foot-Traffic Cleaning \& Surge Monitor Engine}\\[0.5mm]
\small
A smart retail facility records hourly customer foot-traffic counts in a list of raw integer readings \texttt{raw\_readings} representing 8 operational hours:
\begin{itemize}[itemsep=0.5pt,topsep=0.5pt]
  \item \textbf{Sensor Validation \& Cleaning:} Faulty sensor malfunctions record negative counts (e.g., \texttt{-5}, \texttt{-99}). Filter out all negative values into a clean list \texttt{valid\_traffic}.
  \item \textbf{Statistical Summary:} Compute: (1) \texttt{total\_visitors}, (2) \texttt{mean\_hourly}, (3) \texttt{peak\_visitors} (highest single hour), and (4) \texttt{peak\_hour\_idx} (the zero-based index of the peak hour).
  \item \textbf{Surge Detection:} Identify all hours where foot traffic exceeds $1.3 \times \text{mean\_hourly}$. Append their indices to a list \texttt{surge\_hours}.
  \item \textbf{Output Package:} Package the summary into an immutable tuple: \texttt{(total\_visitors, mean\_hourly, peak\_visitors, surge\_hours)}.
\end{itemize}
\end{tcolorbox}

\vspace{0.5mm}
\textbf{Task 1: Identify Sequential Data Structures, Inputs \& Output Tuple} (6 points)
\begin{itemize}[leftmargin=2em]
  \item \textbf{Input List Specification:} \hrulefill\\[3mm]
  \item \textbf{Intermediate Lists for Cleaned Data \& Surges:} \hrulefill\\[3mm]
  \item \textbf{Components of Returned Output Tuple:} \hrulefill
\end{itemize}

\vspace{0.5mm}
\textbf{Task 2: Write Structured Sequential Processing Logic in Ordered Steps} (7 points)\\
Formulate step-by-step logic detailing list initialization, filtering iteration, accumulator/peak scanning, surge detection pass, and tuple assembly.

\begin{answerbox}[height=5.2cm]{Program-Logic Work Area}
\end{answerbox}

\vspace{0.5mm}
\textbf{Task 3: Hand Trace with Given Test Dataset} (6 points)\\
Given \texttt{raw\_readings = [120, -10, 250, 80, -99, 310, 140, 200]}. Trace the cleaned list, mean, peak, and surge hours:

\begin{answerbox}[height=3.5cm]{Calculation Box for Task 3}
\end{answerbox}

\vspace{0.5mm}
\textbf{Task 4: Explain Two Design Decisions} (6 points)
\begin{enumerate}[label=\alph*),leftmargin=2em]
  \item Why is building a new list via \texttt{append()} safer than removing invalid items from \texttt{raw\_readings} in a \texttt{for x in raw\_readings} loop?\\[3.5mm]
  \hrulefill
  \item Why is packaging summary metrics into an immutable \textbf{tuple} advantageous when returning data from a pipeline?\\[3.5mm]
  \hrulefill
\end{enumerate}

\newpage

% ==============================================================================
% PAGE 5: SECTION D, PROBLEM 1
% ==============================================================================

\section*{Section D: Applied Programming in Python \hfill \normalsize\normalfont(30 points)}
\textit{Write complete, syntactically correct Python statements. Utilize lists, tuples, indexing, slicing, and iteration. Do \textbf{not} define custom functions (\texttt{def}) unless explicitly requested. Suggested time: 18 minutes.}

\subsection*{Problem 1: Financial Portfolio Daily Closing Price \& Volatility Analyzer (15 points)}
Write a Python program that processes historical daily closing prices of a security to compute volatility metrics and price momentum.

\textbf{Requirements:}
\begin{enumerate}[leftmargin=2em,label=\arabic*.,itemsep=1pt]
  \item Initialize the list: \texttt{prices = [105.50, 102.00, 108.40, 112.50, 110.00, 115.80, 114.20]}.
  \item Compute the daily price changes: $\Delta_i = \text{prices}[i] - \text{prices}[i-1]$ for $i = 1, \dots, N-1$. Store all changes in a new list \texttt{daily\_changes}.
  \item Find the day index and amount of the \textbf{single largest one-day gain} ($\max \Delta$).
  \item Count the number of \texttt{gain\_days} ($\Delta > 0$) and \texttt{loss\_days} ($\Delta < 0$).
  \item Compute the cumulative percentage return across the entire span: $\text{return} = \frac{\text{prices}[-1] - \text{prices}[0]}{\text{prices}[0]} \times 100$.
  \item Determine the \textbf{maximum drawdown}: the largest peak-to-trough decline between any day $i$ and subsequent day $j$ ($j > i$), or from the running maximum to the current price.
  \item Print a neatly aligned summary report displaying all metrics formatted to 2 decimal places.
\end{enumerate}

The output should follow this format:
\begin{lstlisting}
--- PORTFOLIO PERFORMANCE SUMMARY ---
Analyzed Days: 7 | Initial: $105.50 | Final: $114.20
Cumulative Return: +8.25%
Gain Days: 4 | Loss Days: 2 | Flat Days: 0
Largest Single-Day Surge: +$6.40 (on Day 2)
Max Drawdown from Peak: -$3.50
\end{lstlisting}

\begin{answerbox}[height=13.2cm]{Python Solution for Problem 1}
\end{answerbox}

\newpage

% ==============================================================================
% PAGE 6: SECTION D, PROBLEM 2
% ==============================================================================

\subsection*{Problem 2: Student Examination Gradebook Matrix \& Cohort Analytics (15 points)}
Write a Python program that analyzes academic performance from a 2D matrix representing examination scores for 4 students across 3 core subjects (Subject 1: Python, Subject 2: Linear Algebra, Subject 3: Statistics).

\textbf{Requirements:}
\begin{enumerate}[leftmargin=2em,label=\arabic*.,itemsep=1pt]
  \item Initialize the student cohort matrix:\\
  \texttt{grades = [[85, 90, 78], [72, 68, 80], [95, 92, 98], [60, 75, 70]]}.
  \item Compute the average grade for each individual student (row-wise mean). Store these in a list \texttt{student\_averages}.
  \item Compute the overall cohort average for each subject (column-wise mean). Store these in a list \texttt{subject\_averages}.
  \item Identify the \textbf{top performing student} (student index and highest average).
  \item Identify the \textbf{most challenging subject} (subject column index and lowest subject average).
  \item Count how many students achieved a student average $\ge 80.0$ (Honor Roll students).
  \item Print a clean tabular gradebook displaying each student's ID (\texttt{Student 0}, \texttt{Student 1}, \dots), their 3 individual grades, and their calculated average, followed by the subject averages at the bottom.
\end{enumerate}

The output should follow this format:
\begin{lstlisting}
Student ID    Subj 1   Subj 2   Subj 3   Average 
-------------------------------------------------
Student 0     85       90       78       84.33
Student 1     72       68       80       73.33
Student 2     95       92       98       95.00
Student 3     60       75       70       68.33
-------------------------------------------------
Subject Avg   78.00    81.25    81.50   
Valedictorian: Student 2 (Average: 95.00)
Hardest Subject: Subject 0 (Average: 78.00)
Honor Roll Count: 2
\end{lstlisting}

\begin{answerbox}[height=10.5cm]{Python Solution for Problem 2}
\end{answerbox}

\vfill
\begin{center}
  \rule{0.5\textwidth}{0.4pt}\\[1mm]
  \textbf{\color{primaryblue}--- END OF EXAM PAPER ---}
\end{center}

\end{document}
